\subsection{Таймеры}
\begin{itemize}
  \item \verb|Timer.new(period, fn)| — создаёт \hyperlink{term:timer}{периодический таймер}; \verb|period| — \hyperlink{term:period}{период} в секундах; \verb|fn| вызывается каждый \hyperlink{term:tick}{тик}.
  \item \verb|timer:start()|, \verb|timer:stop()| — запуск/остановка таймера.
  \item \verb|Timer.callLater(delay, fn)| — единичный вызов \verb|fn| через \verb|delay| секунд.
  \item Таймеры и \hyperlink{term:state}{состояния} делайте \hyperlink{term:scope}{глобальными} (без \verb|local|), иначе сборщик мусора их удалит.
\end{itemize}
\textbf{Пример периодического таймера (мигание):}
\begin{minted}{lua}
local ledNumber = 25
local leds = Ledbar.new(ledNumber)
local on = false
function updateBlink()
  local r,g,b = (on and 1 or 0), 0, 0
  for i=0, ledNumber-1 do leds:set(i, r, g, b) end
  on = not on
end
blinkTimer = Timer.new(0.5, updateBlink)  -- 0.5 секунды
blinkTimer:start()
\end{minted}
\textbf{Пример отложенного вызова (через 3 секунды):}
\begin{minted}{lua}
local ledNumber = 25
local leds = Ledbar.new(ledNumber)
Timer.callLater(3, function ()
  for i=0, ledNumber-1 do leds:set(i, 0, 0, 1) end  -- включить синий
end)
\end{minted}
